\usepackage{fontspec}
\usepackage{polyglossia}
\setdefaultlanguage{russian}
\setotherlanguage{english}
\setmainfont{Times New Roman}
\setsansfont{Arial}
\setmonofont{Consolas}

\usepackage[a4paper,margin=2.5cm]{geometry}
\usepackage{setspace}
\onehalfspacing

\usepackage{graphicx}
\usepackage{float}
\usepackage{booktabs}
\usepackage{longtable}
\usepackage{array}
\usepackage{caption}
\usepackage{amsmath,amssymb}
\usepackage{hyperref}
\usepackage{enumitem}
\usepackage[most]{tcolorbox}
\usepackage{titlesec}
\usepackage{fancyhdr}
\usepackage{bookmark}

\hypersetup{
  colorlinks=true,
  linkcolor=black,
  urlcolor=blue,
  citecolor=black
}

\pagestyle{fancy}
\fancyhf{}
\fancyhead[L]{Системный анализ данных для электроэнергетиков}
\fancyhead[R]{\thepage}
\setlength{\headheight}{15pt}

\titleformat{\chapter}[block]
  {\normalfont\bfseries\Large}
  {\thechapter}
  {1em}
  {}

\newcommand{\code}[1]{\texttt{#1}}
\newcommand{\repourl}{https://github.com/Nephalem72/power-data-analysis-handbook}
\newcommand{\repohref}[1]{\href{\repourl}{#1}}
\newcommand{\notebookhref}[2]{\href{\repourl/blob/main/notebooks/#1}{#2}}
\newcommand{\dochref}[2]{\href{\repourl/blob/main/docs/#1}{#2}}

\newcommand{\screenshotplaceholder}[4]{
\begin{figure}[H]
\centering
\begin{tcolorbox}[
  colback=gray!5,
  colframe=black!65,
  width=0.92\textwidth,
  title={#2},
  fonttitle=\bfseries
]
\textbf{Идентификатор рисунка:} #1

\medskip
\textbf{Что должно быть на рисунке:} #3

\medskip
\textbf{Что необходимо сделать:} #4
\end{tcolorbox}
\caption{#2}
\label{#1}
\end{figure}
}

\newcommand{\figureinstruction}[2]{
\begin{tcolorbox}[colback=white,colframe=black!65,width=0.92\textwidth]
\textbf{Подпись к рисунку:} #1

\medskip
\textbf{Текст после рисунка:} #2
\end{tcolorbox}
}

\newcommand{\chaptergoals}[2]{
\begin{tcolorbox}[colback=gray!4,colframe=black!60,title={Цель главы и ожидаемые результаты}]
\textbf{Цель главы.} #1

\medskip
\textbf{После изучения главы обучающийся должен уметь:}
\begin{itemize}[leftmargin=1.25cm]
#2
\end{itemize}
\end{tcolorbox}
}
